\begin{theorem}\label{thm:main,interp}
For $\eps,\delta>0$ with $1-\eps > \delta$ and any fixed integer $k\ge 3$, there exists $\Delta_0(\eps,\delta,k)$ such that for all $\Delta\ge\Delta_0$ the following holds. 
%
Let $G$ be an $n$-vertex, $\Delta$-regular \CUG[f,k,\eps,\delta] instance with $f \ge k^{\eps+\delta-\frac{1}{2}}$.  
Then there is a deterministic algorithm that solves CUG for $G$ in time
\[
n\exp\left(e^{O(\log^2 k)} \log^2\Delta\right)\,.
\]
\end{theorem}
%
\begin{theorem}\label{thm:main,cluster}
For $\eps,\delta>0$ with $1-\eps > \delta$ and any $k$ satisfying $\log k\ge \Delta^{3/2}\log\Delta$, there exists $\Delta_0(\eps,\delta)$ such that for all $\Delta\ge\Delta_0$ the following holds. 
%
Let $G$ be an $n$-vertex, $\Delta$-regular \CUG[f,k,\eps,\delta] instance with $f\ge k^{2\eps+2\delta-1}$.
Then there is a deterministic algorithm that solves CUG for $G$ in time $kn^{O(1)}e^{O(\Delta)}$.
\end{theorem}
%
To illustrate how close Theorem~\ref{thm:main,cluster} gets to an algorithm for usual Unique Games, consider an $n$-vertex, $\Delta$-regular instance of UG and suppose there exists an assignment of value $1-\eps$. 
Let $S$ be an arbitrary set of $\eps n$ vertices, and consider all assignments obtained by relabelling vertices in $S$. 
There are at most $\eps\Delta n$ constraints that could be violated by modifying the labels of vertices in $S$, hence each such assignment has value at least $1-3\eps$.
Thus there are at least $k^{\eps n}$ assignments of value $1-3\eps$, which corresponds to having parameter $f=k^{\eps-1}$ as a CUG instance. 
In summary, if we were permitted to take $f=k^{\eps/3-1}$ in Theorem~\ref{thm:main,cluster}, which is only slightly smaller than the stated $f\ge k^{2\eps+2\delta-1}$, then we would be able to refute the Unique Games Conjecture.
%
The idea at the heart of both Theorem~\ref{thm:main,interp} and Theorem~\ref{thm:main,cluster} is to encode a CUG problem as the problem of approximating the value of a \emph{partition function} $Z(G;w)$ which depends on the instance $G$, and is a polynomial in $w$. 
The partition function is intimately connected to statistical physics, and we use two techniques recently developed for approximating partition functions to prove our two main results. 
Theorem~\ref{thm:main,interp} is proved via \emph{polynomial interpolation}, a method due to Barvinok (see~\cite{Barvinok16} and references therein) and furthered by Patel and Regts~\cite{PR17} who improved the running time in many examples.
 Theorem~\ref{thm:main,cluster} is proved via the \emph{cluster expansion} and methods given in~\cite{BCHPT19}. 
These techniques have recently been developed and applied to the problem of approximating the partition function of the Potts model~\cite{Barvinok16,BDPR18,BCHPT19,LSS19}, random cluster model~\cite{BCHPT19}, and other models from statistical physics~\cite{Barvinok16,HPR19,PR17}.
The main conceptual advances in this paper are demonstrations that these techniques may be adapted to UG instances, cleanly handling the constraints assigned to edges that are not present in the standard Potts model. 
The main technical advance in this paper is a zero-free region for (a generalization of) the ferromagnetic Potts model partition function, that may be of independent interest. See Theorem~\ref{thm:zero-free} below.
%
%

\setcounter{theorem}{5}
